Investigamos a classificação binária do rótulo cardíaco registrado na base pediátrica UCMF, com atenção à possibilidade de vazamento de informação. Auditamos os 21 atributos e excluímos dos preditores motivos de encaminhamento com conteúdo diagnóstico, achados do exame cardiovascular e campos de proveniência incerta. Com 9.361 registros de desenvolvimento, comparamos uma referência majoritária, árvore de decisão, regressão logística e floresta aleatória em validação cruzada aninhada, estratificada e agrupada. Imputação, codificação e escalonamento foram ajustados em cada treino. A floresta alcançou acurácia balanceada de 62,87\% e sensibilidade de 59,29\% no cenário fisiológico. Com quatro atributos antropométricos e demográficos, a regressão logística obteve a maior acurácia balanceada, de 59,74\%. A mudança na ordenação das famílias indica dependência do conjunto de entradas. A cronologia da coleta e a origem do rótulo limitam a interpretação clínica.
